% !TeX root = ../教材题库_全部条目_转写工作稿.tex
% Source: XJ-B1-C06-RV-P268-Y1202
% 第6章 / 章末复习 复习题六 / 教材 p.268、269 / PDF 第275、276页
% OCR未校对：公式与文字均可能有误；未完成数学TeX转写；禁止直接发给学生
% 原题及相邻原解截图（核对用）：../题目原图/XJ-B1-C06-RV-P268-Y1202-01.webp, ../题目原图/XJ-B1-C06-RV-P268-Y1202-02.webp
% 本文件只在转写完成后作为题干使用；OCR 条目尚不是可靠题干。
\begin{bankitem}{复习题六 第16题}{第6章 / 章末复习 复习题六 / 教材 p.268、269 / PDF 第275、276页}{XJ-B1-C06-RV-P268-Y1202}{OCR未校对：公式与文字均可能有误；未完成数学TeX转写；禁止直接发给学生}
\begin{ocrdraft}
[教材p.268 / PDF第275页]
\par
16．一天，在一堂数学课上，老师测量了在场所有学生的身高，计算后发现
\par
男生的平均身高是170cm,女生的平均身高是160cm.其中，张强是最高的，他
\par
的身高是180cm；杜梅是最矮的，她的身高是150cm.
\par
那天有两名学生缺课．但是第二天这两个人均到校上课，老师也测量了他们
\par
的身高，然后重新计算平均身高．有趣的是，女生的平均身高和男生的平均身高
\par
都没变．
\par
下面哪些结论能从这些信息中推断出来？在每个结论的“是”或“否”上面
\par
．两名学生都是女生
\par
是
\par
否
\par
[教材p.269 / PDF第276页]
\par
．一名是男生，另一名是女生
\par
是
\par
否
\par
．两名学生的身高一样
\par
是
\par
否
\par
．所有学生的平均身高没变
\par
是
\par
否
\par
0杜梅仍然是最矮的
\par
是
\par
否
\end{ocrdraft}
\end{bankitem}
